\documentclass[12pt,a4paper]{article}

\input{/home/tabaka/Documents/GitHub/Defs/defs.tex}

\setcounter{zadaniaRozdzial}{3}
\setcounter{section}{0}

\begin{document}
	
	\begin{center}
		\Large Wymagania edukacyjne z matematyki w klasie 2 Technikum na poszczególne oceny.
	\end{center}
	
	Wymagania edukacyjne na poszczególne oceny są stopniowane.
	Uzyskanie danej oceny wymaga spełnienia wymagań przewidzianych
	dla wszystkich ocen niższych, z uwzględnieniem wymagań
	określonych dla danej oceny.
	
	\section{Funkcja liniowa}
	
	\begin{longtable}{|p{16cm}|}
		\hline
		
		\textbf{Wymagania edukacyjne na ocenę dopuszczającą} \\ \hline
		
		\textbullet\ rozpoznaje funkcję liniową na podstawie jej wzoru lub wykresu; \\
		
		\textbullet\ rysuje wykres funkcji liniowej danej wzorem; \\
		
		\textbullet\ określa monotoniczność funkcji liniowej danej wzorem; \\
		
		\textbullet\ oblicza wartość funkcji liniowej dla danego argumentu; \\
		
		\textbullet\ wyznacza miejsce zerowe funkcji liniowej; \\
		
		\textbullet\ odczytuje z wykresu funkcji liniowej zbiór argumentów, dla których funkcja przyjmuje wartości dodatnie lub ujemne; \\
		
		\textbullet\ wyznacza współrzędne punktów przecięcia wykresu funkcji liniowej z osiami układu współrzędnych; \\
		
		\textbullet\ oblicza współczynnik kierunkowy prostej, jeśli dane są współrzędne dwóch punktów należących do tej prostej, a współrzędne tych punktów są liczbami wymiernymi; \\
		
		\textbullet\ rozpoznaje proste równoległe wśród prostych opisanych równaniami kierunkowymi; \\
		
		\textbullet\ rozpoznaje wielkości wprost proporcjonalne. \\ \hline
		
		\textbf{Wymagania edukacyjne na ocenę dostateczną} \\ \hline
		
		\textbullet\ interpretuje współczynniki występujące we wzorze funkcji liniowej; \\
		
		\textbullet\ oblicza argument, dla którego funkcja liniowa przyjmuje daną wartość; \\
		
		\textbullet\ wyznacza algebraicznie zbiory argumentów, dla których funkcja liniowa przyjmuje wartości dodatnie, ujemne, niedodatnie lub nieujemne; \\
		
		\textbullet\ wyznacza równanie prostej przechodzącej przez dwa dane punkty; \\
		
		\textbullet\ przekształca równanie ogólne prostej do postaci kierunkowej i odwrotnie; \\
		
		\textbullet\ sprawdza, czy trzy dane punkty są współliniowe; \\
		
		\textbullet\ rozpoznaje proste prostopadłe wśród prostych opisanych równaniami kierunkowymi; \\
		
		\textbullet\ wyznacza wzór funkcji liniowej, której wykres przechodzi przez dany punkt i jest równoległy do wykresu innej, danej funkcji liniowej; \\
		
		\textbullet\ wyznacza wzór funkcji liniowej, której wykres przechodzi przez dany punkt i jest prostopadły do wykresu innej, danej funkcji liniowej; \\
		
		\textbullet\ rozwiązuje układ równań metodą graficzną; \\
		
		\textbullet\ określa liczbę rozwiązań układu równań liniowych, korzystając z jego interpretacji geometrycznej; \\
		
		\textbullet\ opisuje równaniem zależność między wielkościami wprost proporcjonalnymi. \\ \hline
		
		\textbf{Wymagania edukacyjne na ocenę dobrą} \\ \hline
		
		\textbullet\ wyznacza wartości parametru, dla których funkcja liniowa jest rosnąca, malejąca lub stała; \\
		
		\textbullet\ określa wzajemne położenie prostych na podstawie ich równań; \\
		
		\textbullet\ wyznacza współrzędne wierzchołków wielokąta, gdy dane są równania prostych zawierających jego boki; \\
		
		\textbullet\ analizuje własności funkcji liniowej. \\ \hline
		
		\textbf{Wymagania edukacyjne na ocenę bardzo dobrą} \\ \hline
		
		\textbullet\ oblicza pole figury ograniczonej wykresami funkcji liniowych i osiami układu współrzędnych; \\
		
		\textbullet\ wyznacza wartości parametru, dla których dwie proste są równoległe; \\
		
		\textbullet\ stosuje warunki równoległości i prostopadłości prostych w dowodach dotyczących własności figur geometrycznych; \\
		
		\textbullet\ analizuje własności funkcji liniowej w zależności od wartości współczynników występujących w jej wzorze. \\ \hline
		
		\textbf{Wymagania edukacyjne na ocenę celującą} \\ \hline
		
		\textbullet\ wyprowadza wzór na współczynnik kierunkowy prostej przechodzącej przez dwa punkty; \\
		
		\textbullet\ udowadnia warunek prostopadłości prostych o danych równaniach kierunkowych; \\
		
		\textbullet\ rozwiązuje zadania o znacznym stopniu trudności dotyczące funkcji liniowej. \\ \hline
		
	\end{longtable}
	
	\section{Planimetria}
	
	\begin{longtable}{|p{16cm}|}
		\hline
		
		\textbf{Wymagania edukacyjne na ocenę dopuszczającą} \\ \hline
		
		\textbullet\ rozróżnia trójkąty ostrokątne, prostokątne i rozwartokątne; \\
		
		\textbullet\ stosuje twierdzenie o sumie miar kątów w trójkącie w prostych przypadkach; \\
		
		\textbullet\ uzasadnia przystawanie trójkątów, wykorzystując cechy przystawania, w prostych przypadkach; \\
		
		\textbullet\ zapisuje proporcje długości boków w trójkątach podobnych; \\
		
		\textbullet\ sprawdza, czy dane figury są podobne; \\
		
		\textbullet\ wskazuje w wielokątach odcinki, których długości są proporcjonalne. \\ \hline
		
		\textbf{Wymagania edukacyjne na ocenę dostateczną} \\ \hline
		
		\textbullet\ sprawdza, czy z trzech odcinków o danych długościach można zbudować trójkąt; \\
		
		\textbullet\ uzasadnia podobieństwo trójkątów, wykorzystując cechy podobieństwa, w prostych przypadkach; \\
		
		\textbullet\ wykorzystuje podobieństwo trójkątów do rozwiązywania prostych zadań; \\
		
		\textbullet\ stosuje podobieństwo wielokątów do obliczania długości boków; \\
		
		\textbullet\ stosuje twierdzenie o stosunku pól figur podobnych w prostych zadaniach; \\
		
		\textbullet\ rozwiązuje proste zadania, wykorzystując twierdzenie Talesa. \\ \hline
		
		\textbf{Wymagania edukacyjne na ocenę dobrą} \\ \hline
		
		\textbullet\ przeprowadza dowód twierdzenia o sumie miar kątów w trójkącie; \\
		
		\textbullet\ stosuje cechy przystawania trójkątów do rozwiązywania zadań geometrycznych; \\
		
		\textbullet\ wykorzystuje podobieństwo trójkątów do rozwiązywania problemów praktycznych i zadań geometrycznych; \\
		
		\textbullet\ rozwiązuje zadania dotyczące podobieństwa wielokątów; \\
		
		\textbullet\ rozwiązuje zadania, wykorzystując twierdzenie Talesa; \\
		
		\textbullet\ rozwiązuje zadania wymagające uzasadnienia i dowodzenia z wykorzystaniem przystawania trójkątów. \\ \hline
		
		\textbf{Wymagania edukacyjne na ocenę bardzo dobrą} \\ \hline
		
		\textbullet\ przeprowadza dowód twierdzenia o mierze kąta zewnętrznego trójkąta; \\
		
		\textbullet\ udowadnia, że symetralne boków trójkąta przecinają się w jednym punkcie; \\
		
		\textbullet\ udowadnia, że dwusieczne kątów trójkąta przecinają się w jednym punkcie; \\
		
		\textbullet\ stosuje cechy przystawania trójkątów do rozwiązywania trudniejszych zadań geometrycznych; \\
		
		\textbullet\ wykorzystuje podobieństwo trójkątów do rozwiązywania problemów praktycznych i trudniejszych zadań geometrycznych; \\
		
		\textbullet\ rozwiązuje zadania wymagające uzasadnienia i dowodzenia z wykorzystaniem podobieństwa trójkątów; \\
		
		\textbullet\ rozwiązuje zadania wymagające uzasadnienia i dowodzenia z zastosowaniem twierdzenia Talesa. \\ \hline
		
		\textbf{Wymagania edukacyjne na ocenę celującą} \\ \hline
		
		\textbullet\ przeprowadza dowód twierdzenia Talesa; \\
		
		\textbullet\ rozwiązuje zadania o znacznym stopniu trudności dotyczące przystawania i podobieństwa figur. \\ \hline
		
	\end{longtable}
	
	\section{Funkcja kwadratowa}
	
	\begin{longtable}{|p{16cm}|}
		\hline
		
		\textbf{Wymagania edukacyjne na ocenę dopuszczającą} \\ \hline
		
		\textbullet\ szkicuje wykres funkcji kwadratowej $f(x)=ax^2$, gdzie $a\neq 0$, i odczytuje z wykresu jej własności; \\
		
		\textbullet\ szkicuje wykres funkcji kwadratowej $f(x)=ax^2+q$, gdzie $a\neq 0$, i odczytuje z wykresu jej własności; \\
		
		\textbullet\ szkicuje wykres funkcji kwadratowej $f(x)=a(x-p)^2$, gdzie $a\neq 0$, i odczytuje z wykresu jej własności; \\
		
		\textbullet\ podaje wzór funkcji kwadratowej w postaci ogólnej i kanonicznej; \\
		
		\textbullet\ oblicza współrzędne wierzchołka paraboli i wyznacza równanie jej osi symetrii; \\
		
		\textbullet\ przekształca postać kanoniczną funkcji kwadratowej do postaci ogólnej; \\
		
		\textbullet\ rozwiązuje niepełne równanie kwadratowe metodą wyłączania wspólnego czynnika przed nawias lub z zastosowaniem wzorów skróconego mnożenia; \\
		
		\textbullet\ określa liczbę pierwiastków równania kwadratowego w zależności od znaku wyróżnika; \\
		
		\textbullet\ rozwiązuje równanie kwadratowe w prostych przypadkach, stosując wzory na pierwiastki; \\
		
		\textbullet\ odczytuje miejsca zerowe funkcji kwadratowej zapisanej w postaci iloczynowej; \\
		
		\textbullet\ rozwiązuje nierówność kwadratową w prostych przypadkach. \\ \hline
		
		\textbf{Wymagania edukacyjne na ocenę dostateczną} \\ \hline
		
		\textbullet\ szkicuje wykres funkcji kwadratowej $f(x)=a(x-p)^2+q$, gdzie $a\neq 0$, i odczytuje z wykresu jej własności; \\
		
		\textbullet\ przekształca postać ogólną funkcji kwadratowej do postaci kanonicznej, korzystając ze wzorów na współrzędne wierzchołka paraboli, a następnie szkicuje wykres funkcji i opisuje jej własności; \\
		
		\textbullet\ wyznacza wzór funkcji kwadratowej w postaci ogólnej, gdy dane są współrzędne wierzchołka paraboli i innego punktu jej wykresu; \\
		
		\textbullet\ rozwiązuje równanie kwadratowe, stosując wzory na pierwiastki; \\
		
		\textbullet\ interpretuje geometrycznie rozwiązania równania kwadratowego w zależności od współczynnika $a$ i wyróżnika $\Delta$; \\
		
		\textbullet\ wyznacza algebraicznie współrzędne punktów przecięcia paraboli z osiami układu współrzędnych; \\
		
		\textbullet\ przedstawia trójmian kwadratowy w postaci iloczynowej, jeśli taka postać istnieje w rozważanym zbiorze liczb; \\
		
		\textbullet\ rozwiązuje nierówności kwadratowe; \\
		
		\textbullet\ stosuje pojęcia najmniejszej i największej wartości funkcji oraz wyznacza w prostych przypadkach najmniejszą i największą wartość funkcji kwadratowej w przedziale domkniętym; \\
		
		\textbullet\ analizuje zadanie tekstowe, zapisuje odpowiednie równanie, nierówność lub wzór funkcji kwadratowej opisujące daną zależność i znajduje w prostych przypadkach rozwiązanie spełniające warunki zadania. \\ \hline
		
		\textbf{Wymagania edukacyjne na ocenę dobrą} \\ \hline
		
		\textbullet\ rozwiązuje równania i nierówności kwadratowe w trudniejszych przypadkach; \\
		
		\textbullet\ wykorzystuje postać iloczynową funkcji kwadratowej do rozwiązywania zadań; \\
		
		\textbullet\ stosuje nierówności kwadratowe do wyznaczania dziedziny funkcji zapisanej za pomocą pierwiastka; \\
		
		\textbullet\ wyznacza w trudniejszych przypadkach najmniejszą i największą wartość funkcji kwadratowej w przedziale domkniętym, korzystając z jej własności; \\
		
		\textbullet\ stosuje własności funkcji kwadratowej do rozwiązywania zadań optymalizacyjnych. \\ \hline
		
		\textbf{Wymagania edukacyjne na ocenę bardzo dobrą} \\ \hline
		
		\textbullet\ zaznacza na osi liczbowej sumę, iloczyn i różnicę zbiorów rozwiązań dwóch nierówności kwadratowych; \\
		
		\textbullet\ wykorzystuje postać iloczynową funkcji kwadratowej do rozwiązywania trudniejszych zadań; \\
		
		\textbullet\ stosuje własności funkcji kwadratowej do rozwiązywania trudniejszych zadań optymalizacyjnych; \\
		
		\textbullet\ rozwiązuje zadania tekstowe o podwyższonym stopniu trudności, stosując równania kwadratowe. \\ \hline
		
		\textbf{Wymagania edukacyjne na ocenę celującą} \\ \hline
		
		\textbullet\ wyprowadza wzory na pierwiastki trójmianu kwadratowego; \\
		
		\textbullet\ udowadnia zależności między współczynnikami funkcji kwadratowej w zadaniach o podwyższonym stopniu trudności; \\
		
		\textbullet\ rozwiązuje zadania o znacznym stopniu trudności dotyczące funkcji kwadratowej. \\ \hline
		
	\end{longtable}
	
	\section{Wielomiany}
	
	\begin{longtable}{|p{16cm}|}
		\hline
		
		\textbf{Wymagania edukacyjne na ocenę dopuszczającą} \\ \hline
		
		\textbullet\ podaje przykład wielomianu, określa jego stopień i podaje wartości jego współczynników; \\
		
		\textbullet\ zapisuje wielomian określonego stopnia o danych współczynnikach; \\
		
		\textbullet\ zapisuje wielomian w postaci uporządkowanej; \\
		
		\textbullet\ oblicza wartość wielomianu dla danego argumentu wymiernego; \\
		
		\textbullet\ wyznacza sumę, różnicę i iloczyn wielomianów oraz określa ich stopnie; \\
		
		\textbullet\ określa stopień iloczynu wielomianów bez wykonywania mnożenia; \\
		
		\textbullet\ przekształca wyrażenia algebraiczne, stosując wzory skróconego mnożenia dotyczące wielomianów stopnia drugiego; \\
		
		\textbullet\ rozwiązuje równania wielomianowe, w których wielomian jest zapisany w postaci iloczynu czynników liniowych. \\ \hline
		
		\textbf{Wymagania edukacyjne na ocenę dostateczną} \\ \hline
		
		\textbullet\ oblicza wartość wielomianu dla danego argumentu niewymiernego; \\
		
		\textbullet\ sprawdza, czy dany punkt należy do wykresu wielomianu; \\
		
		\textbullet\ wyznacza współczynnik przy najwyższej potędze oraz wyraz wolny iloczynu wielomianów bez wykonywania mnożenia; \\
		
		\textbullet\ przekształca wyrażenia algebraiczne, stosując wzory skróconego mnożenia dotyczące wielomianów stopnia drugiego; \\
		
		\textbullet\ rozwiązuje równania wielomianowe, w których wielomian jest zapisany w postaci iloczynu czynników stopnia co najwyżej drugiego; \\
		
		\textbullet\ opisuje wielomianem zależności przedstawione w zadaniu i wyznacza dziedzinę tego wielomianu w prostych przypadkach. \\ \hline
		
		\textbf{Wymagania edukacyjne na ocenę dobrą} \\ \hline
		
		\textbullet\ wyznacza współczynniki wielomianu spełniającego podane warunki; \\
		
		\textbullet\ oblicza wartość wielomianu dwóch lub trzech zmiennych dla podanych argumentów; \\
		
		\textbullet\ wyznacza iloczyn wielomianów wielu zmiennych; \\
		
		\textbullet\ rozwiązuje równania wielomianowe w trudniejszych przypadkach, gdy wielomian jest zapisany w postaci iloczynu czynników stopnia co najwyżej drugiego; \\
		
		\textbullet\ sprawdza, czy rozwiązania równania wielomianowego spełniają podane warunki; \\
		
		\textbullet\ opisuje wielomianem pole powierzchni prostopadłościanu i określa dziedzinę tego wielomianu. \\ \hline
		
		\textbf{Wymagania edukacyjne na ocenę bardzo dobrą} \\ \hline
		
		\textbullet\ określa stopień sumy wielomianów w zależności od wartości parametru; \\
		
		\textbullet\ rozwiązuje zadania tekstowe, wykorzystując działania na wielomianach i równania wielomianowe. \\ \hline
		
		\textbf{Wymagania edukacyjne na ocenę celującą} \\ \hline
		
		\textbullet\ rozwiązuje zadania o znacznym stopniu trudności dotyczące wielomianów. \\ \hline
		
	\end{longtable}
	
\end{document}